% TOP_PROBLEM: {"id":30007650,"problem_number":"LOCAL-30007650","title":"Causality behind the Great Gatsby curve","category_id":21,"category":{"id":21,"name":"economics","display_name":"Economics","description":"Open economic theory statements, published research agendas, and proposed scoped specifications; question provenance and historical priority are qualified per record.","slug":"economics","order_index":21,"created_at":"2026-10-08T00:00:00Z"},"status":"provenance_pending","external_url":null,"legacy_ids":[],"published":false,"statement":"In the explicitly specified setting: Does higher inequality itself reduce upward mobility, or do common institutions and statistical features generate much of the observed relationship? The mathematical statement above fixes the target, assumptions and scope of this catalogue formulation.","economics":{"original_id":139,"field":"Inequality and economic mobility","kind":"identification","formalization_basis":"proposed_specification","source_status":"not_verified","priority_status":"not_established","priority_note":"No qualifying problem statement was directly verified, so historical priority cannot be assigned.","earliest_verified_reference":null,"current_reference":{"authors":["Steven N. Durlauf","Andros Kourtellos","Chih Ming Tan"],"title":"The Great Gatsby Curve","year":2022,"url":"https://www.annualreviews.org/content/journals/10.1146/annurev-economics-082321-122703","doi":"10.1146/annurev-economics-082321-122703","locator":"Abstract, paragraphs 1–2","evidence_summary":"The abstract distinguishes mechanical dynamics and multiple behavioral mechanisms, but does not explicitly declare the catalogue causal question open."},"formalization_note":"This is a proposed scoped research specification. No primary passage explicitly stating this entire remaining question as open was verified. The supporting literature may answer important subquestions; this entry must not be advertised as a verified formal open problem."},"statusQualification":"Provenance pending: an explicit published statement of this exact open question was not verified. This is a catalogue-authored candidate specification, not a certified open conjecture. This is a proposed scoped research specification. No primary passage explicitly stating this entire remaining question as open was verified. The supporting literature may answer important subquestions; this entry must not be advertised as a verified formal open problem.","statusReviewedAt":"2026-10-08"}
% FIRST_POSTED: 2026-10-08
% ENTRY_KIND: statement_only
% !TeX program = lualatex
\documentclass[11pt]{article}
\usepackage{fontspec}
\usepackage[a4paper,margin=25mm]{geometry}
\usepackage{amsmath,amssymb,mathtools}
\usepackage{xurl}
\usepackage[hidelinks]{hyperref}
\newcommand{\sref}[2]{\hyperlink{source:#1:#2}{[S#2]}}
\setlength{\emergencystretch}{3em}
\begin{document}
% problemId:30007650
\section*{Causality behind the Great Gatsby curve}\label{30007650}
\subsection{Definitions and mathematical statement}
Fix a collection of regions, a parental cohort, and a child cohort. Let \(F^P_r\) and \(F^C_r\) be the parent and child income distributions in region \(r\). Define inequality \(I_r\) with a predeclared functional, and define persistence \(\beta_r=\operatorname{Cov}(\log Y^P,\log Y^C)/\operatorname{Var}(\log Y^P)\) for positive lifetime incomes with finite log-income second moments and strictly positive parental log-income variance under every compared policy. The observed association between \(I_r\) and \(\beta_r\) is the Great Gatsby relationship, not a causal estimand. Let \(p\) be a specified intervention changing parental inequality, with stated effects on mean income and public resources. Define \(\tau_r(p)=\beta_r(p)-\beta_r(0)\). Identify these policy effects and determine whether they arise through family investments, skills, social influences, political choices, or aspirations. Different interventions producing the same inequality statistic may have different effects; therefore \(\beta\) is not assumed to be a function of inequality alone. Account for changing marginal distributions, measurement error, and regional composition that can create mechanical relationships. The task is a proposed causal specification informed by a review of multiple mechanisms. The reviewed abstract does not itself certify the catalogue's broad causal formulation as an explicitly unsolved problem.

\par\textbf{Formulation and scope.} This is a proposed scoped research specification. No primary passage explicitly stating this entire remaining question as open was verified. The supporting literature may answer important subquestions; this entry must not be advertised as a verified formal open problem.

\par\textbf{Open-question provenance.} An explicit published open-question statement was not verified. Sources below are supporting literature and must not be treated as a first listing.

\par\textbf{Historical priority.} No qualifying problem statement was directly verified, so historical priority cannot be assigned.

\subsection{Short English statement}
In the explicitly specified setting: Does higher inequality itself reduce upward mobility, or do common institutions and statistical features generate much of the observed relationship? The mathematical statement above fixes the target, assumptions and scope of this catalogue formulation.

\subsection{Sources}
\begin{itemize}\raggedright
\item[\textnormal{[S1]}]\hypertarget{source:30007650:1}{} Steven N. Durlauf, Andros Kourtellos, Chih Ming Tan. 2022. The Great Gatsby Curve. Abstract, paragraphs 1–2.
\url{https://www.annualreviews.org/content/journals/10.1146/annurev-economics-082321-122703}.
DOI: 10.1146/annurev-economics-082321-122703.
{\small\textit{Supporting or current literature; see evidence qualification. The abstract distinguishes mechanical dynamics and multiple behavioral mechanisms, but does not explicitly declare the catalogue causal question open.}}
\end{itemize}
\end{document}
